% GENERATED FILE. Edit canonical lessons, then run npm run generate:books.
% canonical-source-hash: fnv1a64:9b2af0fd7afb6b0b
% canonical-lessons: PT-C510-empatar, PT-C510-empate, PT-C510-adversario, PT-C510-derrotar, PT-C510-eliminar

\chapter{Ganhar, empatar, perder}
\label{ch:pt-ganhar-empatar-perder}
\hlchaptermodality{510}

\begin{chapteropening}
\textbf{By the end of this chapter:} \emph{I can talk about match results.}

Complete the last lesson of chapter 510: I can talk about match results.
\end{chapteropening}

\section[empatar]{empatar — to draw, to tie}

\label{lesson:PT-C510-empatar}



\begin{quote}
Before the new one: say the Portuguese for to shoot, then the Portuguese for to head.
\end{quote}

\subsection*{You'll want to know: empatar}
\textbf{empatar} — \textquotedblleft{}to draw, to tie\textquotedblright{}.

\begin{quote}
As duas equipas empataram a zero.
\end{quote}

\begin{quote}
\emph{The two teams drew nil-nil.}
\end{quote}

\emph{Empatar a um} — to draw one-all.

\begin{cousinweb}
\emph{empatar} was borrowed from Spanish \emph{empatar} (EtymDB), which goes back to Italian \emph{patta} \textquotedblleft{}a draw\textquotedblright{}.
\end{cousinweb}

\begin{grammarlens}[title={What you've built}]
One more word for telling, explaining and describing at B1.
\end{grammarlens}

\subsection*{Your turn}
\begin{itemize}
\raggedright
  \item \emph{Say it:} \emph{empatar}
  \item \emph{Say it:} \emph{empatar}, once more
  \item \emph{Say it:} a sentence of your own with \emph{empatar}
  \item \emph{Recall:} say the Portuguese for to shoot, then the Portuguese for to head, then say \emph{empatar} again
\end{itemize}

\begin{tcolorbox}[breakable,colback=teal!4,colframe=teal!35!black,title={Before you move on}]
What does empatar mean? (\textquotedblleft{}To draw, to tie\textquotedblright{}.) Say it once more.
\end{tcolorbox}
\section[o empate]{o empate — draw, tie}

\label{lesson:PT-C510-empate}



\begin{quote}
Before the new one: say the Portuguese for to head, then the Portuguese for to draw.
\end{quote}

\subsection*{You'll want to know: o empate}
\textbf{o empate} — \textquotedblleft{}draw, tie\textquotedblright{}.

\begin{quote}
O jogo acabou num empate.
\end{quote}

\begin{quote}
\emph{The match ended in a draw.}
\end{quote}

The noun of \emph{empatar}.

\begin{cousinweb}
\emph{empate} comes from \emph{empatar} \textquotedblleft{}to draw\textquotedblright{}, which was borrowed from Spanish (EtymDB).
\end{cousinweb}

\begin{grammarlens}[title={What you've built}]
One more word for telling, explaining and describing at B1.
\end{grammarlens}

\subsection*{Your turn}
\begin{itemize}
\raggedright
  \item \emph{Say it:} \emph{o empate}
  \item \emph{Say it:} \emph{o empate}, once more
  \item \emph{Say it:} a sentence of your own with \emph{o empate}
  \item \emph{Recall:} say the Portuguese for to head, then the Portuguese for to draw, then say \emph{o empate} again
\end{itemize}

\begin{tcolorbox}[breakable,colback=teal!4,colframe=teal!35!black,title={Before you move on}]
What does o empate mean? (\textquotedblleft{}Draw, tie\textquotedblright{}.) Say it once more.
\end{tcolorbox}
\section[o adversário]{o adversário — opponent}

\label{lesson:PT-C510-adversario}



\begin{quote}
Before the new one: say the Portuguese for to draw, then the Portuguese for draw.
\end{quote}

\subsection*{You'll want to know: o adversário}
\textbf{o adversário} — \textquotedblleft{}opponent\textquotedblright{}.

\begin{quote}
O nosso adversário de domingo é muito forte.
\end{quote}

\begin{quote}
\emph{Our opponent on Sunday is very strong.}
\end{quote}

Also an adjective: \emph{a equipa adversária}.

\begin{cousinweb}
\emph{adversário} comes from Latin \emph{adversārius} \textquotedblleft{}an opponent\textquotedblright{} (EtymDB), built on \emph{adversus} \textquotedblleft{}turned towards, against\textquotedblright{}.
\end{cousinweb}

\begin{grammarlens}[title={What you've built}]
One more word for telling, explaining and describing at B1.
\end{grammarlens}

\subsection*{Your turn}
\begin{itemize}
\raggedright
  \item \emph{Say it:} \emph{o adversário}
  \item \emph{Say it:} \emph{o adversário}, once more
  \item \emph{Say it:} a sentence of your own with \emph{o adversário}
  \item \emph{Recall:} say the Portuguese for to draw, then the Portuguese for draw, then say \emph{o adversário} again
\end{itemize}

\begin{tcolorbox}[breakable,colback=teal!4,colframe=teal!35!black,title={Before you move on}]
What does o adversário mean? (\textquotedblleft{}Opponent\textquotedblright{}.) Say it once more.
\end{tcolorbox}
\section[derrotar]{derrotar — to defeat, to beat}

\label{lesson:PT-C510-derrotar}



\begin{quote}
Before the new one: say the Portuguese for draw, then the Portuguese for opponent.
\end{quote}

\subsection*{You'll want to know: derrotar}
\textbf{derrotar} — \textquotedblleft{}to defeat, to beat\textquotedblright{}.

\begin{quote}
A equipa da casa derrotou o campeão.
\end{quote}

\begin{quote}
\emph{The home team beat the champions.}
\end{quote}

The noun \emph{a derrota} means defeat.

\begin{cousinweb}
\emph{derrotar} is linked in EtymDB with French \emph{dérouter} \textquotedblleft{}to throw off course, to rout\textquotedblright{}, built on \emph{route} \textquotedblleft{}a road\textquotedblright{}.
\end{cousinweb}

\begin{grammarlens}[title={What you've built}]
One more word for telling, explaining and describing at B1.
\end{grammarlens}

\subsection*{Your turn}
\begin{itemize}
\raggedright
  \item \emph{Say it:} \emph{derrotar}
  \item \emph{Say it:} \emph{derrotar}, once more
  \item \emph{Say it:} a sentence of your own with \emph{derrotar}
  \item \emph{Recall:} say the Portuguese for draw, then the Portuguese for opponent, then say \emph{derrotar} again
\end{itemize}

\begin{tcolorbox}[breakable,colback=teal!4,colframe=teal!35!black,title={Before you move on}]
What does derrotar mean? (\textquotedblleft{}To defeat, to beat\textquotedblright{}.) Say it once more.
\end{tcolorbox}
\section[eliminar]{eliminar — to knock out, to eliminate}

\label{lesson:PT-C510-eliminar}



\begin{quote}
Before the new one: say the Portuguese for opponent, then the Portuguese for to defeat.
\end{quote}

\subsection*{You'll want to know: eliminar}
\textbf{eliminar} — \textquotedblleft{}to knock out, to eliminate\textquotedblright{}.

\begin{quote}
Fomos eliminados nos quartos de final.
\end{quote}

\begin{quote}
\emph{We were knocked out in the quarter-finals.}
\end{quote}

\emph{Ser eliminado} — to be knocked out.

\begin{cousinweb}
\emph{eliminar} is a learned borrowing from Latin \emph{ēlīmināre} \textquotedblleft{}to put out of doors\textquotedblright{} (EtymDB), built on \emph{līmen} \textquotedblleft{}a threshold\textquotedblright{}.
\end{cousinweb}

\begin{grammarlens}[title={What you've built}]
One more word for telling, explaining and describing at B1.
\end{grammarlens}

\subsection*{Your turn}
\begin{itemize}
\raggedright
  \item \emph{Say it:} \emph{eliminar}
  \item \emph{Say it:} \emph{eliminar}, once more
  \item \emph{Say it:} a sentence of your own with \emph{eliminar}
  \item \emph{Recall:} say the Portuguese for opponent, then the Portuguese for to defeat, then say \emph{eliminar} again
\end{itemize}

\begin{tcolorbox}[breakable,colback=teal!4,colframe=teal!35!black,title={Before you move on}]
What does eliminar mean? (\textquotedblleft{}To knock out, to eliminate\textquotedblright{}.) Say it once more.
\end{tcolorbox}
